% Section 3 | Results \& Claims Validation | NIST Quiz \#3 (aligned with CONFIG + JSON artifacts)
\section{Results and Claims Validation}
\label{sec:results-validation}

This section validates the \textbf{14-item trustworthy-AI questionnaire} against on-disk artifacts under \texttt{output/results/}, as indexed by \texttt{test\_cases/CONFIG}. It replaces the legacy applied-cryptography figure checklist (AES/TLS) that previously appeared here.

\subsection{Artifact traceability}
\begin{itemize}
  \item \textbf{Raw completions:} \texttt{output/results/nist\_eval\_latest.json} records \texttt{items[]} with numeric \texttt{id} $1$--$14$, each with \texttt{prompt}, \texttt{response}, \texttt{error}, plus run metadata (\texttt{model}, \texttt{base\_url}, \texttt{temperature}, \texttt{dry\_run}). Claims in Section~\ref{sec:experiments} must be consistent with these strings.
  \item \textbf{Human rubric:} \texttt{output/results/nist\_quiz\_scores.json} holds \texttt{scores} with string keys \texttt{"1"}--\texttt{"14"} and values \texttt{C}, \texttt{P}, or \texttt{N} per \texttt{homework-assignment.pdf}. The LaTeX fragment \texttt{output/results/nist\_rubric\_table.tex} is generated by \texttt{python scripts/emit\_nist\_rubric\_table.py} and included from \texttt{experiment.tex} when building \texttt{src/main.tex}.
  \item \textbf{Threading:} Items \textbf{7} and \textbf{8} form a two-turn sequence in the evaluation driver: item~8 may reference the assistant output from item~7. Validation should confirm the logged \texttt{response} for id~8 reflects that protocol.
\end{itemize}

\subsection{Checks performed}
\begin{enumerate}
  \item \textbf{Completeness:} All fourteen ids appear in \texttt{nist\_eval\_latest.json}; no missing or duplicate \texttt{id} fields.
  \item \textbf{Score coverage:} Every id $1$--$14$ has a rubric entry in \texttt{nist\_quiz\_scores.json} before final submission.
  \item \textbf{Narrative alignment:} Each ``Observed result'' paragraph in \texttt{experiment.tex} (included by \texttt{src/main.tex}) maps to the same item id and \texttt{C}/\texttt{P}/\texttt{N} label as the JSON and rubric table fragment.
  \item \textbf{Dry-run flag:} If \texttt{dry\_run} is \texttt{true}, responses are placeholders; prose should state that a live run is required for empirical claims.
\end{enumerate}

\subsection{Optional demo pipeline (out of scope for rubric proof)}
The scripts under \texttt{examples/python/} (e.g.\ \texttt{generate\_results.py}) produce synthetic \texttt{output/results.json} for workflow demos. They do \textbf{not} replace \texttt{nist\_eval\_latest.json} for Quiz~\#3 grading.
